\documentclass[JEP,XML,SOM,Unicode]{cedram}
\datereceived{2019-05-20}
\dateaccepted{2020-02-24}
\dateepreuves{2020-03-05}

\multlinegap0pt
\newcommand{\loc}{\mathrm{loc}}
\newcommand{\spfrac}[2]{{#1}/{(#2)}}
\newcommand\implique{\mathchoice{\implies}{\Rightarrow}{\Rightarrow}{\Rightarrow}}
\newcommand\ssi{\mathchoice{\iff}{\Leftrightarrow}{\Leftrightarrow}{\Leftrightarrow}}

\DeclareMathOperator{\divv}{div}\let\div\divv

\let\al\alpha
\let \be \beta
\let \de \delta
\let \Om \Omega
\let \ph \varphi
\let \si \sigma

\def \NN {\mathbb N}
\def \RR {\mathbb R}
\def \Ss {\mathbb S}
\def \TT {\mathbb T}
\def \CC {\mathbb C}

\theoremstyle{plain}
\newtheorem{Thm}{Theorem}[section]
\newtheorem{Pro}[Thm]{Proposition}
\newtheorem{Lem}[Thm]{Lemma}

\theoremstyle{definition}
\newtheorem{Rem}[Thm]{Remark}

\begin{document}
\frontmatter
\title[Homogenization of linear transport equations. A new approach]{Homogenization of\\ linear transport equations.\\ A new approach}

\author[\initial{M.} \lastname{Briane}]{\firstname{Marc} \lastname{Briane}}
\address{Univ Rennes, INSA Rennes, CNRS, IRMAR - UMR 6625\\
F-35000 Rennes, France}
\email{mbriane@insa-rennes.fr}
\urladdr{http://briane.perso.math.cnrs.fr/}

\begin{abstract}
The paper is devoted to a new approach of the homogenization of linear transport equations induced by a uniformly bounded sequence of vector fields $b_\varepsilon(x)$, the solutions of which $u_\varepsilon(t,x)$ agree at $t=0$ with a bounded sequence of $L^p_\mathrm{loc}(\mathbb{R}^N)$ for some $p\in(1,\infty)$. Assuming that the sequence $b_\varepsilon\cdot\nabla w_\varepsilon^1$ is compact in $L^q_\mathrm{loc}(\mathbb{R}^N)$ ($q$ conjugate of $p$) for some gradient field $\nabla w_\varepsilon^1$ bounded in $L^N_\mathrm{loc}(\mathbb{R}^N)^N$, and that there exists a uniformly bounded sequence $\sigma_\varepsilon>0$ such that $\sigma_\varepsilon\,b_\varepsilon$ is divergence free if $N\!=\!2$ or is a cross product of $(N\!-\!1)$ bounded gradients in $L^N_\mathrm{loc}(\mathbb{R}^N)^N$ if $N\!\geq\!3$, we prove that the sequence $\sigma_\varepsilon\,u_\varepsilon$ converges weakly to a solution to a linear transport equation. It turns out that the compactness of $b_\varepsilon\cdot\nabla w_\varepsilon^1$ is a substitute to the ergodic assumption of the classical two-dimensional periodic case, and allows us to deal with non-periodic vector fields in any dimension. The homogenization result is illustrated by various and general examples.
\end{abstract}

\keywords{Homogenization, transport equation, dynamic flow, rectification}

\alttitle{Homogénéisation d'équations de transport linéaires. Une nouvelle approche}

\altkeywords{Homogénéisation, équation de transport, flot dynamique, redressement}

\begin{altabstract}
Cet article propose une nouvelle approche de l'homogénéisation des équations de transport linéaires induites par une suite uniformément bornée de champs de vecteurs $b_\varepsilon(x)$ et dont les solutions $u_\varepsilon(t,x)$ coïncident en $t=0$ avec une suite bornée de $L^p_\mathrm{loc}(\mathbb{R}^N)$ pour un certain $p\in(1,\infty)$. En supposant que la suite $b_\varepsilon\cdot\nabla w_\varepsilon^1$ est compacte dans $L^q_\mathrm{loc}(\mathbb{R}^N)$ ($q$ exposant conjugué de $p$) pour un champ de gradients $\nabla w_\varepsilon^1$ borné dans $L^N_\mathrm{loc}(\mathbb{R}^N)^N$ et qu'il existe une suite uniformément bornée $\sigma_\varepsilon>0$ telle que $\sigma_\varepsilon\,b_\varepsilon$ est à divergence nulle si $N\!=\!2$ ou est un produit vectoriel de $(N\!-\!1)$ gradients bornés dans $L^N_\mathrm{loc}(\mathbb{R}^N)^N$ si $N\!\geq\!3$, on montre que la suite $\sigma_\varepsilon\,u_\varepsilon$ converge faiblement vers une solution d'une équation de transport. Il s'avère que la compacité de $b_\varepsilon\cdot\nabla w_\varepsilon^1$ remplace la condition d'ergodicité du cas périodique bidimensionnel classique et permet de traiter des champs de vecteurs non périodiques en toute dimension. Le résultat d'homogénéisation est illustré par différents exemples généraux.
\end{altabstract}

\subjclass{35B27, 35F05, 37C10}
\maketitle
\tableofcontents
\mainmatter

\section{Introduction}
In this paper we study the homogenization of the sequence of linear transport equations indexed by $\varepsilon>0$,
\begin{equation}\label{eque}
\begin{cases}
\dfrac{\partial u_\varepsilon}{\partial t}-b_\varepsilon(x)\cdot\nabla_x u_\varepsilon=0 & \mbox{in }(0,T)\times\RR^N,
\\[5pt]
u_\varepsilon(0,\cdot)=u_\varepsilon^0 & \mbox{in }\RR^N.
\end{cases}
\end{equation}
where $N\geq 2$, $T>0$ and $p\in[1,\infty]$ with conjugate exponent $q$.
Equation \eqref{eque} is associated with the flow $Y_\varepsilon(t,x)$ defined by
\begin{equation}\label{Yebei}
\begin{cases}
\dfrac{\partial Y_\varepsilon(t,x)}{\partial t}=b_\varepsilon(Y_\varepsilon(t,x)), & t\in\RR
\\[5pt]
Y_\varepsilon(0,x)=x\in\RR^N.
\end{cases}
\end{equation}
Using the DiPerna-Lions transport theory \cite[Cor.\,II.1]{DiLi}, if for instance $b_\varepsilon$ is a vector field in $L^\infty(\RR^N)^N\cap W^{1,q}_\loc(\RR^N)^N$ with bounded divergence and the initial condition~$u_\varepsilon^0$ is in $L^p(\RR^N)$, then there exists a unique solution $u_\varepsilon(t,x)$ to equation \eqref{eque} in $L^\infty(0,T;L^p(\RR^N))$.

Tartar \cite{Tar} showed that the homogenization of first-order hyperbolic equations may lead to nonlocal effective equations with memory effects.
In this framework Amirat, Hamdache and Ziani obtained memory effects for the homogenization of the transport equation \eqref{eque}, first in dimension two when $b_\varepsilon\cdot\nabla_x u_\varepsilon=b_\varepsilon(t,x_2)\,\partial_{x_1}u_\varepsilon$ \cite{AHZ1}, and in a more general setting \cite{AHZ2,AHZ3} they derived an integral representation of the memory term of the homogenized equation using Young's measures and a Radon transform.

In the present paper, rather than determining the ``closure'' set of the homogenized equations of \eqref{eque} as in \cite{AHZ1,AHZ2,AHZ3}, we provide some sufficient conditions under which the homogenized equation of \eqref{eque} remains a transport equation, namely the nature of equation \eqref{eque} is preserved through homogenization.
Assuming that $b_\varepsilon(x)=b(x/\varepsilon)$ is a divergence free periodically oscillating vector field, Brenier~\cite{Bre} showed thanks to the ergodic theorem that the solution $u_\varepsilon$ of \eqref{eque} converges. This result was improved in dimension two by Hou and Xin \cite{HoXi} who proved thanks to a two-scale convergence approach that the homogenization of \eqref{eque} leads to the averaged transport equation with $\int_{\TT^2}b(y)\,dy$ ($\TT^N$ denotes the $N$-dimensional torus). Golse \cite[Th.\,8]{Gol1} extended these results to the locally periodic case $b_\varepsilon(x)=b(x,x/\varepsilon)$ with $\div_yb(x,\cdot)=0$, assuming some ergodic property involving the flow associated with $-\,b(x,\cdot)$. In \cite{Gol1,Gol2} and in the generalizations \cite{DGL,DuGo} with an application to Physics \cite{DEG}, Golse et al.\ also studied the perturbed differential system satisfied by the pair $(I_\varepsilon,J_\varepsilon)$ in $\RR^N\times\TT^N$:
\begin{equation}\label{IeJe}
\begin{cases}
\dfrac{d I_\varepsilon}{dt}=\varepsilon f(I_\varepsilon,J_\varepsilon), & t\in\RR
\\[8pt]
\dfrac{d J_\varepsilon}{dt}=\omega(I_\varepsilon)+\varepsilon g(I_\varepsilon,J_\varepsilon),
\end{cases}
\end{equation}
with $f(I,\cdot)$, $g(I,\cdot)$ periodic, which after an $\varepsilon$-rescaling is associated with a Liouville partial differential equation of type \eqref{eque} but more complicated. Assuming among others that $\omega$ satisfies the Kolmogorov non-degeneracy type ergodic condition
\begin{equation}\label{ergK}
\mbox{meas}\left(\big\{I\in\RR^N,\ |I|\leq R:|\omega(I)\cdot\xi|<\al\big\}\right)\;\xrightarrow[\al\to 0]\; 0\quad\mbox{uniformly in }\xi\in\Ss^{N-1},
\end{equation}
or some variant, Golse et al.\ obtained an error estimate between the solution to system~\eqref{IeJe} and the averaged system with $\int_{\TT^N}f(I,y)\,dy$, as well as the velocity averaging (homogenization) of the associated Liouville partial differential equation. They also extended this result to a non-periodic framework.

Returning to the periodic case $b_\varepsilon(x)=b(x/\varepsilon)$, Tassa \cite{Tas} replaced in dimension two the divergence free of $b$ by the existence of a periodic positive regular function $\sigma$ such that
\begin{equation}\label{divb}
\div\,(\sigma b)=0\quad\mbox{in }\RR^2,
\end{equation}
\ie $\sigma$ is an invariant measure for $b$ by the Liouville theorem.
The main assumption of the periodic framework $b_\varepsilon(x)=b(x/\varepsilon)$ of \cite{Bre,HoXi,Tas} is the ergodicity of the flow associated with $b$ (see, \eg \cite[Lect.\,1]{Sin}, or \cite[Chap.\,II, \S\,5]{ReSi}), namely any periodic invariant function by the flow is constant, or equivalently, for any periodic regular function $v$,
\begin{equation}\label{erg}
b\cdot\nabla v=0\;\;\mbox{in }\RR^2\;\implique\;\nabla v=0\;\;\mbox{in }\RR^2,
\end{equation}
together with $b\neq 0$ in $\RR^2$. By virtue of the Kolmogorov theorem (see, \eg \cite[Lect.\,11]{Sin} or \cite[\S 2]{Tas}) in dimension two with $b\neq 0$, condition \eqref{erg} is equivalent to the ergodic assumption
\begin{equation}\label{birr}
\int_{\TT^2}b_1(y)\,\sigma(y)\,dy,\int_{\TT^2}b_2(y)\,\sigma(y)\,dy\quad\mbox{are rationally independent},
\end{equation}
which is equivalent to the irrationality of the rotation number.
In the locally periodic case $b_\varepsilon(x)=b(x,x/\varepsilon)$ of \cite[\S 8]{Gol1} the ergodicity assumption states that for a.e. $x\in\RR^N$, the fluctuation of the vector field $b(x,\cdot)$ around its average value is ergodic with respect to the flow associated with $-\,b(x,\cdot)$.

In the present paper, besides the closure results of \cite{AHZ1,AHZ2,AHZ3} and the ergodic approaches of \cite{Bre,HoXi,Gol1,Gol2,Tas}, we propose a new approach which holds both in a non-periodic framework and in any dimension, assuming that the vector field $b_\varepsilon$ satisfies a non-ergodic condition which preserves the nature of equation \eqref{eque} through homogenization.
More precisely, the ergodic assumption \eqref{erg} or \eqref{birr} of the periodic framework
is now replaced by the existence of a sequence $w_\varepsilon^1$ in $C^1(\RR^N)$ and $q\in(1,\infty)$ such that
\begin{equation}\label{beDw1e}
0< b_\varepsilon\cdot\nabla w_\varepsilon^1\;\to\;\theta_0>0\quad\mbox{strongly in }L^q_\loc(\RR^N),
\end{equation}
which is equivalent in the periodic case to the existence of a periodic gradient $\nabla w$ satisfying
\begin{equation}\label{bDw1}
b\cdot\nabla w=1\quad\mbox{in }\RR^N.
\end{equation}
Moreover, the invariant measure $\sigma$ of the periodic case is replaced by a sequence $\sigma_\varepsilon$ satisfying $0<c^{-1}<\sigma_\varepsilon<c$ for some constant $c>1$, and (see Remark~\ref{rem.rec} for an equivalent expression)
\begin{equation}\label{siebei}
\div\,(\sigma_\varepsilon\,b_\varepsilon)=0\;\;\mbox{if }N=2\quad\mbox{and}\quad \sigma_\varepsilon\,b_\varepsilon=\nabla w_\varepsilon^2\times\cdots\times \nabla w_\varepsilon^N\;\;\mbox{if }N\geq 3.
\end{equation}
The case where $\sigma_\varepsilon\,b_\varepsilon$ is only divergence free in dimension $N\geq 3$ remains open.
In this way the vector field $b_\varepsilon$ is naturally associated with the vector field $W_\varepsilon:=(w_\varepsilon^1,\dots, w_\varepsilon^N)$ which induces a global rectification of the field $b_\varepsilon$ in the direction~$e_1$ (see Remark~\ref{rem.rec}). Then, assuming in addition to \eqref{beDw1e}, \eqref{siebei} that $W_\varepsilon$ is uniformly proper (see condition \eqref{upWe} below) and converges both in $C^0_\loc(\RR^N)^N$ and weakly in $W^{1,N}_\loc(\RR^N)^N$, we prove (see Theorem~\ref{thm.hom}) that, up to a subsequence, $\sigma_\varepsilon\,u_\varepsilon$ converges weakly in $L^\infty(0,T;L^p(\RR^N))$ to a solution $v$ to the transport equation
\begin{equation}\label{eqvi}
\begin{cases}
\dfrac{\partial v}{\partial t}-\xi_0\cdot\nabla_x\Bigl(\dfrac{v}{\sigma_0}\Bigr)=0 & \mbox{in }(0,T)\times\RR^N
\\[5pt]
v(0,\cdot)=v^0 & \mbox{in }\RR^N,
\end{cases}
\end{equation}
where $\sigma_0$ is the weak-$\ast$ limit of $\sigma_\varepsilon$ in $L^\infty(\RR^N)$, $\xi_0$ is the weak limit of $\sigma_\varepsilon\,b_\varepsilon$ in $L^{\spfrac{N}{N-1}}_\loc(\RR^N)^N$ and $v^0$ the weak limit of $\sigma_\varepsilon\,u_\varepsilon^0$ in $L^p(\RR^N)$. Moreover, if $\sigma_\varepsilon$ converges strongly to $\sigma_0$ in $L^1_\loc(\RR^N)$ (see Remark~\ref{rem.vsiu}) or $u_\varepsilon^0$ converges strongly to $u^0$ in $L^p_\loc(\RR^N)$, then up to a subsequence $u_\varepsilon$ converges weakly in $L^\infty(0,T;L^p(\RR^N))$ to a solution $u$ to the transport equation
\begin{equation}\label{equ1}
\begin{cases}
\dfrac{\partial u}{\partial t}-\dfrac{\xi_0}{\sigma_0}\cdot\nabla_x u=0 & \mbox{in }(0,T)\times\RR^N
\\[5pt]
u(0,\cdot)=u^0 & \mbox{in }\RR^N.
\end{cases}
\end{equation}
The convergence of $u_\varepsilon$ also turns out to be strong in $L^\infty(0,T;L^2_\loc(\RR^N))$ if $u_\varepsilon^0$ converges strongly to $u^0$ in $L^p_\loc(\RR^N)$ with $p>2$ (see the second part of Theorem~\ref{thm.hom}).

The compactness condition \eqref{beDw1e} is the main assumption of Theorem~\ref{thm.hom}. It is equivalent to the compactness of the product $\sigma_\varepsilon\det(DW_\varepsilon)$ which is connected to the vector field $b_\varepsilon$ by \eqref{siebei}. The examples of Section~\ref{s.exa} show that this condition may be satisfied in quite general situations.

Section~\ref{s.res} is devoted to the statement of the main result and to its proof.
Section~\ref{s.exa} deals by three applications of Theorem~\ref{thm.hom}.
In Section~\ref{ss.exa1} we study the case of a diffeomorphism $W_\varepsilon$ on $\RR^2$ such that $\det(DW_\varepsilon)$ is compact in $L^p_\loc(\RR^2)$ for some $q\in\nobreak (1,\infty)$. In Section~\ref{ss.per} we extend the periodic case of \cite{Bre,HoXi,Tas} with $b_\varepsilon(x)=b(x/\varepsilon)$ and the periodic case of \cite[\S 4]{Bri} on the asymptotic of the flow associated with $b$, in~the light of Theorem~\ref{thm.hom} with a periodically oscillating function $\sigma_\varepsilon(x)=\sigma(x/\varepsilon)$ (see~ Proposition~\ref{pro.per}). In Section~\ref{ss.flow} we consider the case of a diffeomorphism $W_\varepsilon$ which agrees at a fixed time $t$ to a flow $X_\varepsilon(t,\cdot)$ associated with a suitable vector field $a_\varepsilon$ (see~Proposition~\ref{pro.ae}). In this general setting assumption~\eqref{beDw1e} holds simply when $\div\,a_\varepsilon$ is compact in $L^q_\loc(\RR^N)$ for some $q\in(1,\infty)$.

\skpt
\subsubsection*{Notations}
\begin{itemize}
\item $\left(e_1,\dots,e_N\right)$ denotes the canonical basis of $\RR^N$.
\item $\cdot$ denotes the scalar product in $\RR^N$ and $|\cdot|$ the associated norm.
\item $I_N$ is the unit matrix of $\RR^{N\times N}$, and $R_\perp$ is the clockwise $90^\circ$ rotation matrix in $\RR^{2\times 2}$.
\item For $M\in\RR^{N\times N}$, $M^T$ denotes the transpose of $M$.
\item $Y_N:=\left[0,1\right)^N$, and $\langle f\rangle$ denotes the average-value of a function $f\in L^1(Y_N)$.

\item For any open set $\Om$ of $\RR^N$ and $k\in\NN\cup\{\infty\}$, $C^k_c(\Om)$, respectively $C^k_b(\Om)$, denotes the space of the $C^k$ functions with compact support in $\Om$, respectively bounded in $\Om$.
\item For $k\in\NN\cup\{\infty\}$ and $p\in[1,\infty]$, $C^k_\sharp(Y_N)$ denotes the space of the $Y_N$-periodic functions in $C_k(\RR^N)$, and $L^p_\sharp(Y_N)$ denotes the space of the $Y_N$-periodic functions in $L^p_\loc(\RR^N)$ (\ie in $L^p(K)$ for any compact set $K$ of $\RR^N$).
\item For $u\in L^1_\loc(\RR^N)$ and $U=(U_j)_{1\leq j\leq d}\in L^1_\loc(\RR^N)^N$.
\[
\nabla_x u:=\left(\partial_{x_1}u,\dots,\partial_{x_N}u\right)\quad\mbox{and}\quad DU:=\big[\partial_{x_i}U_j\big]_{1\leq i,j\leq d}.
\]
The matrix-valued function $DU$ stands for the transposition of the Jacobian matrix of the vector field $U$.
\item For $\xi^1,\dots,\xi^N$ in $\RR^N$, the cross product $\xi^2\times\cdots\times \xi^N$ is defined by
\begin{equation}\label{crpr}
\xi^1\cdot(\xi^2\times\dots\times \xi^N)=\det\,(\xi^1,\xi^2,\dots,\xi^N)\quad\mbox{for }\xi^1\in\RR^N,
\end{equation}
where $\det$ is the determinant with respect to the canonical basis $(e_1,\dots,e_N)$.
\item $o_\varepsilon$ denotes a term which tends to zero as $\varepsilon\to 0$.
\item $C$ denotes a constant which may vary from line to line.
\end{itemize}

\subsubsection*{Acknowledgements} The author is very grateful to the unknown referees for several comments which have improved the presentation of the paper and enlarged the bibliography.

\section{The main result}\label{s.res}
Let $W_\varepsilon=(w_\varepsilon^1,\dots,w_\varepsilon^N)$, $\varepsilon>0$, be a sequence of vector fields in $C^1(\RR^N)^N$
which is {\em uniformly proper}, \ie for any compact set $K$ of $\RR^N$ there exists a compact subset $K'$ of $\RR^N$ satisfying\enlargethispage{-\baselineskip}%
\begin{equation}\label{upWe}
W_\varepsilon^{-1}(K)\subset K'\quad\mbox{for any small enough }\varepsilon>0,
\end{equation}
and let $W\in C^1(\RR^N)^N$ be such that
\begin{equation}\label{cWe}
W_\varepsilon\to W\;\;\mbox{in }C^0_\loc(\RR^N)^N\quad\mbox{and}\quad W_\varepsilon\rightharpoonup W\;\;\mbox{in }W^{1,N}_\loc(\RR^N)^N.
\end{equation}
Let $b_\varepsilon$ be a vector field in $C^0_b(\RR^N)^N\cap W^{1,q}_\loc(\RR^N)^N$ with bounded divergence and let~$\sigma_\varepsilon$ be a positive function in $C^0(\RR^N)\cap W^{1,q}_\loc(\RR^N)$ satisfying for some constant $c>1$,
\begin{equation}\label{siebe}
c^{-1}\leq\sigma_\varepsilon\leq c\quad\mbox{and}\quad
\sigma_\varepsilon\,b_\varepsilon=\begin{cases}
R_\perp\nabla w_\varepsilon^2 & \mbox{if }N=2
\\[3pt]
\nabla w_\varepsilon^2\times\cdots\times \nabla w_\varepsilon^N & \mbox{if }N\geq 3,
\end{cases}
\quad \mbox{in }\RR^N.
\end{equation}
Also assume that for $p\in(1,\infty)$ with conjugate exponent $q$, there exists a positive function $\theta_0$ in $C^0(\RR^N)$ such that
\begin{equation}\label{cthe}
\theta_\varepsilon:=b_\varepsilon\cdot\nabla w_\varepsilon^1>0\quad\mbox{in }\RR^N \qquad\mbox{and}\qquad \theta_\varepsilon\to\theta_0>0\quad\mbox{strongly in }L^q_\loc(\RR^N).
\end{equation}
\pagebreak[2]%
Finally, assume:
\begin{itemize}
\item either that there exists a constant $B>0$ such that
\begin{equation}\label{beB}
|\div\,b_\varepsilon|\leq B\quad\mbox{a.e. in }\RR^N,
\end{equation}
\item or the regularity condition
\begin{equation}\label{rsiebe}
b_\varepsilon\in C^1_b(\RR^N)^N,\quad \sigma_\varepsilon\in C^1(\RR^N)\quad\mbox{and}\quad u_\varepsilon^0\in C^1(\RR^N).
\end{equation}
\end{itemize}
\begin{Rem}\label{rem.rec}
The definition \eqref{siebe} of $b_\varepsilon$ can be also written for any dimension $N\geq 2$ as the existence of $(N-1)$ gradients $\nabla w_\varepsilon^2,\dots,\nabla w_\varepsilon^N$ satisfying
\begin{equation}\label{siebe2}
\forall\,\xi\in\RR^N,\quad \si_\varepsilon\,b_\varepsilon\cdot\xi=\det\left(\xi,\nabla w_\varepsilon^2,\dots,\nabla w_\varepsilon^N\right).
\end{equation}
In dimension $N\geq 3$ this is exactly the definition of the cross product $\nabla w_\varepsilon^2\times\cdots\times\nabla w_\varepsilon^N$ (see~\eqref{crpr}). In dimension $N=2$ this means exactly that $\si_\varepsilon\,b_\varepsilon=R_\perp\nabla w_\varepsilon^2$, which is equivalent to
\begin{equation}\label{divsiebe}
\div\,(\sigma_\varepsilon\,b_\varepsilon)=0\;\;\mbox{in }\RR^2.
\end{equation}
However, in dimension $N\geq 3$ condition \eqref{siebe} is stronger than $\sigma_\varepsilon\,b_\varepsilon$ divergence free.

The definition \eqref{siebe} of $b_\varepsilon$ and the definition \eqref{cthe} of $\theta_\varepsilon$ are equivalent to the global rectification of the field $b_\varepsilon$ by the diffeomorphism $W_\varepsilon$
\begin{equation}\label{DWebe}
DW_\varepsilon^T\,b_\varepsilon=\theta_\varepsilon\,e_1\quad\mbox{in }\RR^N,
\end{equation}
in the direction $e_1$ with the compact range $\theta_\varepsilon$.
\end{Rem}

Then, we have the following homogenization result.
\begin{Thm}\label{thm.hom}
Let $T>0$, let $p\in(1,\infty)$ and let $u_\varepsilon^0$ be a bounded sequence in $L^p(\RR^N)$, $N\geq 2$. Assume that conditions \eqref{upWe} to \eqref{cthe} together with \eqref{beB} or \eqref{rsiebe} hold true.
Let $u_\varepsilon$ be the solution to the transport equation \eqref{eque} and set $v_\varepsilon:=\sigma_\varepsilon\,u_\varepsilon$.
Then, up to a subsequence $v_\varepsilon$ converges weakly in $L^\infty(0,T;L^p(\RR^N))$ to a solution $v$ to the transport equation
\begin{equation}\label{eqv}
\begin{cases}
\dfrac{\partial v}{\partial t}-\xi_0\cdot\nabla_x\Bigl(\dfrac{v}{\sigma_0}\Bigr)=0 & \mbox{in }(0,T)\times\RR^N
\\
v(0,\cdot)=v^0 & \mbox{in }\RR^N,
\end{cases}
\end{equation}
where (${\rm Cof}$ denotes the cofactors matrix)
\begin{gather}\label{xi0cofDW}
\xi_0={\rm Cof}\,(DW)\,e_1\in C^0(\RR^N)^N,
\\
\label{xisiv0}
\begin{cases}
\sigma_\varepsilon\,b_\varepsilon\rightharpoonup \xi_0\;\;\mbox{in }L^{\spfrac{N}{N-1}}_\loc(\RR^N)^N,\\
\sigma_\varepsilon\rightharpoonup\sigma_0\;\;\mbox{in }L^\infty(\RR^N)\,\ast,\\
\sigma_\varepsilon\,u_\varepsilon^0\rightharpoonup v^0\;\;\mbox{in }L^p(\RR^N).
\end{cases}
\end{gather}
Moreover, if in addition $b_\varepsilon\in W^{1,p/(p-2)}_\loc(\RR^N)^N$ with $p>2$ and the sequence $u_\varepsilon^0$ converges strongly to $u^0\in L^p_\loc(\RR^N)$ with
\[
\sigma_0\in W^{1,\infty}(\RR)\quad\text{and}\quad\xi_0\in L^\infty(\RR^N)^N\cap W^{1,p/(p-2)}_\loc(\RR^N)^N,
\]
then $u_\varepsilon$ converges strongly in $L^\infty(0,T;L^2_\loc(\RR^N))$ to the solution $u$ to the transport equation
\begin{equation}\label{equ}
\begin{cases}
\dfrac{\partial u}{\partial t}-\dfrac{\xi_0}{\sigma_0}\cdot\nabla_x u=0 & \mbox{in }(0,T)\times\RR^N
\\[5pt]
u(0,\cdot)=u^0 & \mbox{in }\RR^N.
\end{cases}
\end{equation}
\end{Thm}
\begin{Rem}\label{rem.uni}
If in Theorem~\ref{thm.hom} we assume in addition that $\sigma_0$ is in $W^{1,\infty}(\RR^N)$ and~$\xi_0$ belongs to $L^\infty(\RR^N)^N\cap W^{1,q}_\loc(\RR^N)^N$, then by virtue of \cite[Cor.\,II.1]{DiLi} there exists a unique solution $v$ to the transport equation \eqref{eqv}.
\end{Rem}
\begin{Rem}\label{rem.vsiu}
In addition to the conditions \eqref{upWe} to \eqref{cthe} assume that $\sigma_\varepsilon$ converges strongly in $L^1_\loc(\RR^N)$ to $\sigma_0\in W^{1,q}_\loc(\RR^N)$. Then, we have $v=\sigma_0\,u$ and $v^0=\sigma_0\,u^0$ where~$u^0$ is the weak limit of $u_\varepsilon^0$ in $L^p(\RR^N)$, which implies that equation \eqref{eqv} is equivalent to equation \eqref{equ}. Therefore, $u_\varepsilon$ converges weakly in $L^\infty(0,T;L^p(\RR^N))$ to a solution $u$ to the transport equation \eqref{equ}.
\end{Rem}
To prove Theorem~\ref{thm.hom} we need the following $L^p$-estimate.
\begin{Lem}\label{lem.Lpeue}
Let $b_\varepsilon\in L^\infty(\RR^N)^N\cap W^{1,q}_\loc(\RR^N)^N$ with bounded divergence be such that
\begin{itemize}
\item either estimate \eqref{beB} holds true,
\item or both conditions \eqref{siebe} and \eqref{rsiebe} hold true.
\end{itemize}
Then, there exists a constant $C>0$ such that for any $u_\varepsilon^0\in L^p(\RR^N)$ with $p\in[1,\infty)$, the solution $u_\varepsilon$ to equation~\eqref{eque} satisfies the estimate
\begin{equation}\label{Lpeue}
\|u_\varepsilon(t,\cdot)\|_{L^p(\RR^N)}\leq C\,\|u_\varepsilon^0\|_{L^p(\RR^N)}\quad\mbox{for a.e. }t\in(0,T),
\end{equation}
\end{Lem}

\begin{proof}[Proof of Theorem~\ref{thm.hom}]
First of all, note that by \eqref{siebe} and \eqref{cthe} we have
\begin{equation}\label{detDWe}
\det (DW_\varepsilon)=\sigma_\varepsilon\,\theta_\varepsilon>0\quad\mbox{in }\RR^N.
\end{equation}
This combined with property \eqref{upWe} and Hadamard-Caccioppoli's theorem \cite{Cac} (or~Hada\-mard-Lévy's theorem) implies that $W_\varepsilon$ is a $C^1$-diffeomorphism on $\RR^N$. Moreover, since by \eqref{detDWe} $\det(DW_\varepsilon)$ is positive and by \eqref{cWe} $W_\varepsilon$ converges weakly in $W^{1,N}_\loc(\RR^N)^N$, by~virtue of Müller's theorem \cite{Mul} $\det(DW_\varepsilon)$ weakly converges to $\det(DW)$ in $L^1_\loc(\RR^N)$. Hence, passing to the limit in \eqref{detDWe} together with the strong convergence~\eqref{cthe} of $\theta_\varepsilon$, the weak convergence \eqref{xisiv0} of $\sigma_\varepsilon$ and the boundedness \eqref{siebe} of $\sigma_\varepsilon$ we get that
\begin{equation}\label{detDW}
\det (DW)=\sigma_0\,\theta_0\geq c^{-1}\,\theta_0>0\quad\mbox{a.e. in }\RR^N,
\end{equation}
which taking into account the continuity of $DW$ and $\theta_0$ implies that $\det (DW)>0$ in $\RR^N$.
Moreover, again by the uniform character of \eqref{upWe} $W$ is a proper mapping.
Therefore, $W$ is also a $C^1$-diffeomorphism on $\RR^N$.
The weak formulation of equation \eqref{eque} is that for any function $\phi\in C^1_c([0,T)\times\RR^N)$,
\begin{equation}\label{weque}
\int_0^T\!\!\int_{\RR^N}u_\varepsilon\,\dfrac{\partial\phi}{\partial t}\,dx\,dt+\int_{\RR^N}u_\varepsilon^0(x)\,\phi(0,x)\,dx=\int_0^T\!\!\int_{\RR^N}u_\varepsilon\,\div\left(\phi\,b_\varepsilon\right)dx\,dt.
\end{equation}
Using a density argument with $\sigma_\varepsilon\in W^{1,q}_\loc(\RR^N)$, we can replace the test function $\phi$ by $\sigma_\varepsilon\,\ph$ for any $\ph\in C^1_c([0,T)\times\RR^N)$. This combined with the divergence free of $\sigma_\varepsilon\,b_\varepsilon$ leads us to the new formulation
\begin{multline}\label{eqsieue}
\int_0^T\!\!\int_{\RR^N}\,\sigma_\varepsilon\,u_\varepsilon\,\dfrac{\partial\ph}{\partial t}\,dx\,dt+\int_{\RR^N}\sigma_\varepsilon(x)\,u_\varepsilon^0(x)\,\ph(0,x)\,dx\\[-8pt]
=\int_0^T\!\!\int_{\RR^N}u_\varepsilon\,\sigma_\varepsilon\,b_\varepsilon\cdot\nabla_x\ph\,dx\,dt.
\end{multline}
We pass easily to the limit in the left hand-side of \eqref{eqsieue}. The delicate point comes from the right-hand side of \eqref{eqsieue}.

By the $L^p$-estimate \eqref{Lpeue} of Lemma~\ref{lem.Lpeue} combined with the uniform boundedness of $\sigma_\varepsilon$ in \eqref{siebe} there exists a subsequence, still denoted by $\varepsilon$, such that $v_\varepsilon=\sigma_\varepsilon\,u_\varepsilon$ converges weakly to some function $v$ in $L^\infty(0,T;L^p(\RR^N))$.

Let $\psi\in C^1_c([0,T)\times\RR^N)$ the support of which is contained in some compact set $[t_0,t_1]\times K$ of $[0,T)\times\RR^N$, and define
\begin{equation}\label{phe}
\ph_\varepsilon(t,x):=\psi(t, W_\varepsilon(x))\quad\mbox{for }(t,x)\in(0,T)\times\RR^N,
\end{equation}
so that $\nabla_x\ph_\varepsilon(t,x)=DW_\varepsilon(x)\nabla_y\psi(t,y)$.
Hence, making the change of variables $y=W_\varepsilon(x)$ and using \eqref{DWebe} we deduce that
\begin{multline}\label{vebe}
\int_0^T\!\!\int_{\RR^N}v_\varepsilon(t,x)\,b_\varepsilon(x)\cdot\nabla_x\ph_\varepsilon(t,x)\,dx\,dt\\
=\int_0^T\!\!\int_{W_\varepsilon^{-1}(K)}v_\varepsilon(t,x)\,b_\varepsilon(x)\cdot\nabla_x\ph_\varepsilon(t,x)\,dx\,dt
\\
=\int_0^T\!\!\int_Kv_\varepsilon(t,W_\varepsilon^{-1}(y))\,\theta_\varepsilon(W_\varepsilon^{-1}(y))\,e_1\cdot\nabla_y\psi(t,y)\,\det(DW_\varepsilon^{-1})(y)\,dy\,dt.
\end{multline}
First, using successively the Hölder inequality combined with the $L^p$-estimate \eqref{Lpeue}, the inclusion \eqref{upWe} and the $L^q$-strong convergence \eqref{cthe} of $\theta_\varepsilon$, we have
\begin{multline*}
\biggl|\,\int_0^T\!\!\int_{K}v_\varepsilon(t,W_\varepsilon^{-1}(y))\,(\theta_\varepsilon-\theta_0)(W_\varepsilon^{-1}(y))\,e_1\cdot\nabla_y\psi(t,y)\,\det(DW_\varepsilon^{-1})(y)\,dy\,dt\,\biggr|
\\
\leq C_\psi \int_0^T\!\left(\int_{K}\!\big|v_\varepsilon(t,W_\varepsilon^{-1}(y))\big|^p\det(DW_\varepsilon^{-1})(y)\,dy\right)^{1/ p}
\\
\shoveright{\times\left(\int_{K}\!\big|(\theta_\varepsilon-\theta_0)(W_\varepsilon^{-1}(y))\big|^q\det(DW_\varepsilon^{-1})(y)\,dy\right)^{1/ q}\!\!dt}
\\
\leq C_\psi\int_0^T\|v_\varepsilon(t,\cdot)\|_{L^p(K')}\|\theta_\varepsilon-\theta_0\|_{L^q(K')}\,dt=o_\varepsilon,
\end{multline*}
which implies that
\begin{multline*}
 \int_0^T\!\!\int_{K}v_\varepsilon(t,W_\varepsilon^{-1}(y))\,\theta_\varepsilon(W_\varepsilon^{-1}(y))\,e_1\cdot\nabla_y\psi(t,y)\,\det(DW_\varepsilon^{-1})(y)\,dy\,dt
\\
=\int_0^T\!\!\int_{K}v_\varepsilon(t,W_\varepsilon^{-1}(y))\,\theta_0(W_\varepsilon^{-1}(y))\,e_1\cdot\nabla_y\psi(t,y)\,\det(DW_\varepsilon^{-1})(y)\,dy\,dt+o_\varepsilon.
\end{multline*}
Next, by the uniform convergence~\eqref{cWe}
\[
\nabla_y\psi(t,W_\varepsilon(x))\to\nabla_y\psi(t,W(x))\quad\mbox{in }C^0_\loc([0,T]\times\RR^N).
\]
Then, making the inverse change of variables $x=W_\varepsilon^{-1}(y)$ together with \eqref{upWe} and using the weak convergence of $v_\varepsilon$ to $v$ in $L^\infty(0,T;L^p(\RR^N))$, we have
\begin{multline*}
 \int_0^T\!\!\int_{K}v_\varepsilon(t,W_\varepsilon^{-1}(y))\,\theta_0(W_\varepsilon^{-1}(y))\,e_1\cdot\nabla_y\psi(t,y)\,\det(DW_\varepsilon^{-1})(y)\,dy\,dt
\\
= \int_0^T\!\!\int_{K'}v_\varepsilon(t,x)\,\theta_0(x)\,e_1\cdot\nabla_y\psi(t,W_\varepsilon(x))\,dx\,dt
\\
=\int_0^T\!\!\int_{K'}v(t,x)\,\theta_0(x)\,e_1\cdot\nabla_y\psi(t,W(x))\,dx\,dt+o_\varepsilon.
\end{multline*}
Let $\ph\in C^1_c([0,T)\times\RR^N)$ and define similarly to \eqref{phe}
\[
\ph(t,x):=\psi(t, W(x))\quad\mbox{for }(t,x)\in[0,T)\times\RR^N,
\]
so that $\nabla_x\ph(t,x)=DW(x)\nabla_y\psi(t,y)$.
Therefore, passing to the limit in \eqref{vebe} we obtain that
\begin{multline}\label{cvebe}
\int_0^T\!\!\int_{\RR^N}v_\varepsilon(t,x)\,b_\varepsilon(x)\cdot\nabla_x\ph_\varepsilon(t,x)\,dx\,dt
\\
=\int_0^T\!\!\int_{\RR^N}v(t,x)\,\theta_0(x)\,\big(DW(x)^T\big)^{-1}e_1\cdot\nabla_x\ph(t,x)\,dx\,dt+o_\varepsilon.
\end{multline}
On the other hand, using \eqref{DWebe}, \eqref{siebe} and the Murat-Tartar div-curl lemma in $L^{\spfrac{N}{N-1}}$-$L^N$ (see, \eg \cite[Th.\,2]{Mur}) with convergences \eqref{cWe}, \eqref{cthe}, \eqref{xisiv0} we get that
\begin{equation}\label{cDWesiebe}
DW_\varepsilon^T(\sigma_\varepsilon\,b_\varepsilon)=\sigma_\varepsilon\,\theta_\varepsilon\,e_1\rightharpoonup DW^T\xi_0=\sigma_0\,\theta_0\,e_1\quad\mbox{weakly in }L^1_\loc(\RR^N).
\end{equation}
This combined with \eqref{detDW} yields equality \eqref{xi0cofDW}.
Convergences \eqref{cvebe} and \eqref{cDWesiebe} imply that
\[
\int_0^T\!\!\int_{\RR^N}v_\varepsilon\,b_\varepsilon\cdot\nabla_x\ph_\varepsilon\,dx\,dt\;\xrightarrow[\varepsilon\to 0]\;
\int_0^T\!\!\int_{\RR^N}\dfrac{v}{\sigma_0}\,\xi_0\cdot\nabla_x\ph\,dx\,dt.
\]
Finally, passing to the limit in formula \eqref{eqsieue} with $\ph_\varepsilon$, it follows that for any $\ph\in C^1_c([0,T)\times\RR^N)$,
\[
\int_0^T\!\!\int_{\RR^N}\,v\, \dfrac{\partial\ph}{\partial t}\,dx\,dt+\int_{\RR^N}v^0(x)\,\ph(0,x)\,dx
=\int_0^T\!\!\int_{\RR^N} \dfrac{v}{\sigma_0}\,\xi_0\cdot\nabla_x\ph\,dx\,dt,
\]
which taking into account that $\xi_0$ is divergence free yields the weak formulation of the desired limit equation \eqref{eqv}.
This concludes the proof of the first part of Theorem~\ref{thm.hom}.

Now, assume in addition that $b_\varepsilon\in W^{1,p/(p-2)}_\loc(\RR^N)^N$ with $p>2$ and $u_\varepsilon^0$ converges strongly to $u^0$ in $L^p(\RR^N)$ with
\[
\sigma_0\in W^{1,\infty}(\RR^N)\quad\text{and}\quad \xi_0\in L^\infty(\RR^N)^N\cap W^{1,p/(p-2)}_\loc(\RR^N)^N.
\]
By \cite[Th.\,II.3 and Corollary\,II.1]{DiLi} $u_\varepsilon^2$ is the unique solution to the equation \eqref{eque} with initial condition $(u_\varepsilon^0)^2$, or equivalently, for any $\phi\in C^1_c([0,T)\times\RR^N)$,
\[
\int_0^T\!\!\int_{\RR^N}u_\varepsilon^2\, \dfrac{\partial\phi}{\partial t}\,dx\,dt+\int_{\RR^N}(u_\varepsilon^0)^2(x)\,\phi(0,x)\,dx=\int_0^T\!\!\int_{\RR^N}u_\varepsilon^2\,\div\left(\phi\,b_\varepsilon\right)dx\,dt,
\]
Replacing $u_\varepsilon$ by $u_\varepsilon^2$ in the first part of Theorem~\ref{thm.hom} and using the strong convergence of $u_\varepsilon^0$ we get that the sequence $\sigma_\varepsilon\,u_\varepsilon^2$ converges weakly in $L^\infty(0,T;L^{p/2}(\RR^N))$ to the solution $w$ to the transport equation
\begin{equation}\label{eqw}
\begin{cases}
\dfrac{\partial w}{\partial t}-\xi_0\cdot\nabla_x\Bigl(\dfrac{w}{\sigma_0}\Bigr)= \dfrac{\partial w}{\partial t}-\dfrac{\xi_0}{\sigma_0}\cdot\nabla_x w+ \dfrac{\xi_0\cdot\nabla\sigma_0}{\sigma_0^2}\,w=0
& \mbox{in }(0,T)\times\RR^N
\\[8pt]
w(0,\cdot)=\sigma_0\,(u^0)^2 & \mbox{in }\RR^N.
\end{cases}
\end{equation}
Note that by virtue of \cite[Cor.\,II.1]{DiLi} the solution $w$ to equation \eqref{eqw} is unique due to the regularities $\sigma_0\in W^{1,\infty}(\RR^N)$, $\xi_0\in L^\infty(\RR^N)^N\cap W^{1,p/(p-2)}_\loc(\RR^N)^N$ with divergence free. Moreover, again by \cite[Th.\,II.3 and Corollary\,II.1]{DiLi} $v^2$ is the unique solution to the equation induced by \eqref{eqv}
\[
\begin{cases}
\dfrac{\partial(v^2)}{\partial t}-\dfrac{\xi_0}{\sigma_0}\cdot\nabla_x(v^2)+2\, \dfrac{\xi_0\cdot\nabla\sigma_0}{\sigma_0}\,v^2=0 & \mbox{in }(0,T)\times\RR^N
\\[5pt]
v^2(0,\cdot)=(\sigma_0\, u^0)^2 & \mbox{in }\RR^N,
\end{cases}
\]
or equivalently, for any $\phi\in C^1_c([0,T)\times\RR^N)$,
\begin{multline*}
 \int_0^T\!\!\int_{\RR^N}v^2\, \dfrac{\partial\phi}{\partial t}\,dx\,dt+\int_{\RR^N}(\sigma_0\, u^0)^2(x)\,\phi(0,x)\,dx
\\
=\int_0^T\!\!\int_{\RR^N}v^2\,\div\Bigl(\phi\, \dfrac{\xi_0}{\sigma_0}\Bigr)dx\,dt+\int_0^T\!\!\int_{\RR^N}2\,v^2\, \dfrac{\xi_0\cdot\nabla\sigma_0}{\sigma_0^2}\,\phi\,dx\,dt.
\end{multline*}
Replacing the test function $\phi$ by $\ph/\sigma_0$ by a density argument, it follows that for any function $\ph\in C^1_c([0,T)\times\RR^N)$,
\begin{multline*}
 \int_0^T\!\!\int_{\RR^N} \dfrac{v^2}{\sigma_0}\, \dfrac{\partial\ph}{\partial t}\,dx\,dt+\int_{\RR^N}\sigma_0(x)\,(u^0)^2(x)\,\ph(0,x)\,dx
\\
=\int_0^T\!\!\int_{\RR^N}v^2\,\div\Bigl(\ph\, \dfrac{\xi_0}{\sigma_0^2}\Bigr)dx\,dt+\int_0^T\!\!\int_{\RR^N}2\,v^2\, \dfrac{\xi_0\cdot\nabla\sigma_0}{\sigma_0^3}\,\ph\,dx\,dt
\\
=\int_0^T\!\!\int_{\RR^N} \dfrac{v^2}{\sigma_0}\,\div\Bigl(\ph\, \dfrac{\xi_0}{\sigma_0}\Bigr)dx\,dt+\int_0^T\!\!\int_{\RR^N} \dfrac{v^2}{\sigma_0}\, \dfrac{\xi_0\cdot\nabla\sigma_0}{\sigma_0^2}\,\ph\,dx\,dt,
\end{multline*}
which shows that $v^2/\sigma_0$ is also a solution to equation \eqref{eqw}. By uniqueness we thus get that $w=v^2/\sigma_0$.
Similarly, the solution $u$ to equation \eqref{equ} agrees with $v/\sigma_0$.
Finally, using these two equalities we have for any compact set $K$ of $\RR^N$,
\begin{multline*}
 \int_0^T\!\!\int_{K}\sigma_\varepsilon(u_\varepsilon-u)^2\,dx\,dt=\int_0^T\!\!\int_{K}(\sigma_\varepsilon\,u_\varepsilon^2-2\,\sigma_\varepsilon\,u_\varepsilon\,u+\sigma_\varepsilon\,u^2)\,dx\,dt
\\[-8pt]
\xrightarrow[\varepsilon\to 0]\;\int_0^T\!\!\int_{K}(w-2\,v\,u+\sigma_0\,u^2)\,dx\,dt=0,
\end{multline*}
which concludes the proof of Theorem~\ref{thm.hom}.
\end{proof}

\begin{proof}[Proof of Lemma~\ref{lem.Lpeue}]
If the uniform boundedness \eqref{beB} of $\div\,b_\varepsilon$ is satisfied, then using the estimate (17) of \cite[Prop.\,II.1]{DiLi} for the solution to the regularized equation of~\eqref{eque} and the lower semi-continuity of the $L^p$-norm ($p<\infty)$ we get estimate \eqref{Lpeue}.

Otherwise, assume that conditions \eqref{siebe} and \eqref{rsiebe} hold true.
Using the regularity of the data the proof is based on an explicit expression of the solution to equation~\eqref{eque} from the flow $Y_\varepsilon$ associated with the vector field $b_\varepsilon$ by
\begin{equation}\label{Yebe}
\begin{cases}
\dfrac{\partial Y_\varepsilon(t,x)}{\partial t}=b_\varepsilon(Y_\varepsilon(t,x)), & t\in\RR
\\[5pt]
Y_\varepsilon(0,x)=x\in\RR^N.
\end{cases}
\end{equation}
Let $u_\varepsilon^0$ be a function in $C^1(\RR^N)\cap L^p(\RR^N)$.
It is classical that the regular solution $u_\varepsilon$ to the transport equation \eqref{eque} is given by
\begin{equation}\label{ueXe}
u_\varepsilon(t,x)=u_\varepsilon^0(Y_\varepsilon(t,x))\quad\mbox{for }(t,x)\in[0,T]\times\RR^N.
\end{equation}
Let $t\in[0,T]$. Making the change of variables combined with the semi-group property of the flow
\[
y=Y_\varepsilon(t,x)\;\ssi\;x=Y_\varepsilon(-t,y),
\]
we get that
\begin{equation}\label{Lpeue01}
\int_{\RR^N}\big|u_\varepsilon^0(Y_\varepsilon(t,x))\big|^p\,dx=\int_{\RR^N}\big|u_\varepsilon^0(y)\big|^p\,\big|\det(D_y Y_\varepsilon(-t,y))\big|\,dy.
\end{equation}
Moreover, by \eqref{Yebe} and the Liouville formula we have for any $(\tau,y)\in\RR\times\RR^N$,
\[
\det(D_y Y_\varepsilon(\tau,y))=\exp\biggl(\int_0^{\tau}(\div\,b_\varepsilon)(Y_\varepsilon(s,y))\,ds\biggr).
\]
However, since by \eqref{siebe} $\sigma_\varepsilon\,b_\varepsilon$ is divergence free, we have
\begin{multline*}
 \int_0^{\tau}(\div\,b_\varepsilon)(Y_\varepsilon(s,y))\,ds=-\int_0^{\tau}\Bigl(\dfrac{\nabla\sigma_\varepsilon\cdot b_\varepsilon}{\sigma_\varepsilon}\Bigr)(Y_\varepsilon(s,y))\,ds
\\
=-\int_0^{\tau} \dfrac{\partial}{\partial s}\big(\ln \sigma_\varepsilon(Y_\varepsilon(s,y))\big)\,ds=\ln\Bigl(\dfrac{\sigma_\varepsilon(y)}{\sigma_\varepsilon(Y_\varepsilon(\tau,y))}\Bigr).
\end{multline*}
This combined with the boundedness of $\sigma_\varepsilon$ in condition \eqref{siebe} implies that
\[
\forall\,(\tau,y)\in\RR\times\RR^N,\quad 0<\det(D_y Y_\varepsilon(\tau,y))= \dfrac{\sigma_\varepsilon(y)}{\sigma_\varepsilon(Y_\varepsilon(\tau,y))}\leq c^2.
\]
Hence, we deduce from \eqref{Lpeue01} that
\[
\int_{\RR^N}|u_\varepsilon(x)|^p\,dx=\int_{\RR^N}\big|u_\varepsilon^0(Y_\varepsilon(t,x))\big|^p\,dx\leq c^2\int_{\RR^N}\big|u_\varepsilon^0(y)\big|^p\,dy,
\]
which yields the desired estimate \eqref{Lpeue}. This concludes the proof of Lemma~\ref{lem.Lpeue}.
\end{proof}


\section{Examples}\label{s.exa}
The purpose of this section is to illustrate the homogenization of the transport equation \eqref{eque} by various oscillating fields $b_\varepsilon$ which satisfy the assumptions of Theorem~\ref{thm.hom}. It means giving examples of diffeomorphism $W_\varepsilon$ on $\RR^N$ satisfying the rectification \eqref{DWebe} of the vector field $b_\varepsilon$ where the sequence $\theta_\varepsilon>0$ is compact in $L^q_\loc(\RR^N)$ for some $q\in(1,\infty)$.

\subsection{First example}\label{ss.exa1}
Let $\al_\varepsilon,\al\in C^1(\RR)$ be such that for some constant $c>0$,
\begin{equation}\label{ae}
\al_\varepsilon\to\al\;\;\mbox{in }C^0_\loc(\RR),\quad \al_\varepsilon'\geq c\;\;\mbox{in }\RR,\quad\al_\varepsilon'\to\al'\;\;\mbox{in }L^2_\loc(\RR),
\end{equation}
and let $\be_\varepsilon,\be\in C^1(\RR)$ be such that for some constant $C>0$,
\begin{equation}\label{bee}
\be_\varepsilon\to\be\;\;\mbox{in }C^0_\loc(\RR),\quad |\be_\varepsilon|\leq C\;\;\mbox{in }\RR,\quad \be_\varepsilon'\mbox{ is bounded in }L^\infty_\loc(\RR),
\end{equation}
Consider the vector field $W_\varepsilon\in C^1(\RR^N)^N$ defined by
\begin{multline}\label{Wecj}
W_\varepsilon(x)\\
:=\begin{pmatrix}\al_\varepsilon(x_1)\,\exp\big\{\be_\varepsilon(\al_\varepsilon(x_1)\al_\varepsilon(x_2))\big\},\al_\varepsilon(x_2)\,\exp\big\{\!-\!\be_\varepsilon(\al_\varepsilon(x_1)\al_\varepsilon(x_2))\big\}\end{pmatrix},\quad
x\in\RR^2,
\end{multline}
which is based on the characterization of the holomorphic mappings on $\CC^2$ with constant Jacobian \cite{Nis}.
The gradient of $W_\varepsilon$ is given by
\[
\begin{cases}
\nabla w_\varepsilon^1(x)=\exp\big\{\be_\varepsilon(\al_\varepsilon(x_1)\al_\varepsilon(x_2))\big\}\!\begin{pmatrix}
\al_\varepsilon'(x_1)\big(1+\al_\varepsilon(x_1)\al_\varepsilon(x_2)\be_\varepsilon'(\al_\varepsilon(x_1)\al_\varepsilon(x_2))\big)
\\
\al_\varepsilon'(x_2)\al_\varepsilon^2(x_1)\be_\varepsilon'(\al_\varepsilon(x_1)\al_\varepsilon(x_2)\big)
\end{pmatrix}
\\
\nabla w_\varepsilon^2(x)\\[-8pt]
\hspace*{.5cm}=\exp\big\{\!-\!\be_\varepsilon(\al_\varepsilon(x_1)\al_\varepsilon(x_2))\big\}\!\begin{pmatrix}
-\al_\varepsilon'(x_1)\al_\varepsilon^2(x_2)\be_\varepsilon'(\al_\varepsilon(x_1)\al_\varepsilon(x_2))
\\
\al_\varepsilon'(x_2)\big(1-\al_\varepsilon(x_1)\,\al_\varepsilon(x_2)\be_\varepsilon'(\al_\varepsilon(x_1)\al_\varepsilon(x_2))\big)
\end{pmatrix}.
\end{cases}
\]
Also define $b_\varepsilon:=R_\perp\nabla w_\varepsilon^2$ and $\sigma_\varepsilon:=1$, so that conditions \eqref{siebe} and \eqref{beB} are fulfilled.

By \eqref{ae} and \eqref{bee} we have
\begin{gather*}
 W_\varepsilon(x)\xrightarrow{C^0_\loc(\RR^2)}W(x)\!:=\!\bigl(\al(x_1)\,\exp\big\{\be(\al(x_1)\al(x_2))\big\},\al(x_2)\,\exp\big\{\!-\!\be(\al(x_1)\al(x_2))\big\}\bigr),
\\
W_\varepsilon\rightharpoonup W \quad\mbox{in }H^1_\loc(\RR^2),
\end{gather*}
so that conditions \eqref{cWe} are satisfied, and
\begin{equation}\label{jacWe}
b_\varepsilon\cdot\nabla w_\varepsilon^1(x)=\det(DW_\varepsilon)(x)=\al_\varepsilon'(x_1)\,\al_\varepsilon'(x_2)\to \al'(x_1)\,\al'(x_2)\quad\mbox{in }L^2_\loc(\RR^2),
\end{equation}
so that condition \eqref{cthe} is satisfied with $p=2$.
Moreover, since by \eqref{ae}
\[
\forall\,t\in\RR,\quad |\al_\varepsilon(t)-\al_\varepsilon(0)|\geq c\,|t|,
\]
the sequence $\al_\varepsilon(0)$ converges, and $\be_\varepsilon$ is uniformly bounded in $\RR$, condition \eqref{upWe} holds for $W_\varepsilon$.

Note that the oscillations of the drift $b_\varepsilon$ in equation \eqref{eque} are only due to the oscillations of the sequence $\be'_\varepsilon$ which does not appear in the convergence \eqref{jacWe} of the Jacobian.

\subsection{The periodic case}\label{ss.per}
This section extends the periodic framework of \cite{Bre,HoXi,Tas}, \cite[Th.\,8]{Gol1}, and \cite[Cor.\,4.4]{Bri}.

Let $W=(w^1,\dots,w^N)$ be a vector field in $C^2(\RR^N)^N$, and let $M$ be a matrix in $\RR^{N\times N}$ such that
\begin{equation}\label{siW}
\big(x\mapsto W(x)-Mx\big)\in C^1_\sharp(Y_N)^N\quad\mbox{and}\quad \sigma:=\det(DW)>0\;\;\mbox{in }\RR^N.
\end{equation}
Consider the periodic vector field $b\in C^1_\sharp(Y^N)^N$ defined by
\begin{equation}\label{sib}
\sigma\,b:=\begin{cases}
R_\perp\nabla w^2 & \mbox{if }N=2
\\[3pt]
\nabla w^2\times\cdots\times \nabla w^N & \mbox{if }N\geq 3.
\end{cases}
\end{equation}

We have the following result.
\begin{Pro}\label{pro.per}
Let $u_\varepsilon^0\in C^1(\RR^N)$ be a bounded sequence in $L^p(\RR^N)$ with $p\in(1,\infty)$.
Assume that conditions \eqref{siW} and \eqref{sib} hold true.
Then, the vector fields $W_\varepsilon$ and $b_\varepsilon$ defined by
\begin{equation}\label{Webe}
W_\varepsilon(x):=\varepsilon\,W\left(\sfrac{x}{\varepsilon}\right)\quad\mbox{and}\quad b_\varepsilon(x):=b\left(\sfrac{x}{\varepsilon}\right)\quad\mbox{for }x\in\RR^N,
\end{equation}
satisfy the assumptions of Theorem~\ref{thm.hom}.

Moreover, for any sequence $u_\varepsilon^0$ in $L^p(\RR^N)$ such that $\sigma(x/\varepsilon)\,u_\varepsilon^0$ converges weakly to $v^0$ in $L^p(\RR^N)$, the solution $u_\varepsilon$ to equation \eqref{eque} is such that $\sigma(x/\varepsilon)\,u_\varepsilon$ converges weakly in $L^\infty(0,T;L^p(\RR^N))$ to the solution $v$ to the equation \eqref{eqv} with $\sigma_0=\langle\sigma\rangle$ and $\xi_0=\langle\sigma\,b\rangle$.
\end{Pro}

\begin{proof}[Proof of Proposition~\ref{pro.per}]
By the quasi-affinity of the determinant (see, \eg \cite[\S 4.3.2]{Dac}) and by \eqref{siW} we have
\[
\det(M)=\det\,\langle DW\rangle=\big\langle \det(DW)\big\rangle >0,
\]
and by \eqref{Webe} there exists a constant $C>0$ such that
\begin{equation}\label{eWeM}
\forall\,x\in\RR^N,\quad |W_\varepsilon(x)-Mx| \leq C\varepsilon,
\end{equation}
which implies condition \eqref{upWe}.
Moreover, estimate \eqref{eWeM} and the uniform bounded of $DW_\varepsilon$ imply easily the convergences \eqref{cWe} with the limit $W(x):=Mx$.

On the other hand, the definitions \eqref{siW} of $W$, $\sigma$ and the definition \eqref{sib} of $b$ show clearly that condition \eqref{siebe} and the regularity \eqref{rsiebe} hold true.
Moreover, we have
\[
\theta:=b\cdot\nabla w^1= \dfrac{\det(DW)}{\sigma}=1\quad\mbox{in }\RR^N,
\]
which implies \eqref{cthe} since $\theta_\varepsilon(x)=\theta(x/\varepsilon)=1$.

Moreover, let $u_\varepsilon^0$ be a sequence in $L^p(\RR^N)$ such that $\sigma(x/\varepsilon)\,u_\varepsilon$ converges weakly to~$v^0$ in $L^p(\RR^N)$. By virtue of Theorem~\ref{thm.hom} combined with Remark~\ref{rem.uni} and using the weak limit of a periodically oscillating sequence, the sequence $\sigma(x/\varepsilon)\,u_\varepsilon$ converges weakly in $L^p(\RR^N)$ to the solution $v$ to the equation \eqref{eqv} with $\sigma_0=\langle\sigma\rangle$ and $\xi_0=\langle\sigma\,b\rangle$.
The proof of Proposition~\ref{pro.per} is now complete.
\end{proof}

\subsection{The dynamic flow case}\label{ss.flow}
In this section we construct a sequence $W_\varepsilon$ from a dynamic flow associated with a suitable but quite general sequence of vector fields $a_\varepsilon$.

Let $a_\varepsilon,a$ be vector fields in $C^1(\RR^N)^N$ such that
\begin{equation}\label{cae}
a_\varepsilon \to a\;\;\mbox{in }C^0_\loc(\RR^N)^N,\quad a_\varepsilon\rightharpoonup a\;\;\mbox{in }W^{1,\infty}_\loc(\RR^N)^N\ast,
\end{equation}
and for some constant $A>0$,
\begin{equation}\label{aeA}
|a_\varepsilon|+|\div\,a_\varepsilon|\leq A\quad\mbox{in }\RR^N.
\end{equation}
Also assume that there exists $q\in(1,\infty)$ such that
\begin{equation}\label{divae}
\div\,a_\varepsilon\to \div\,a\quad\mbox{in }L^q_\loc(\RR^N).
\end{equation}
Consider the dynamic flow $X_\varepsilon$ associated with the vector field $a_\varepsilon$ defined by
\begin{equation}\label{Xeae}
\begin{cases}
\dfrac{\partial X_\varepsilon(t,x)}{\partial t}=a_\varepsilon(X_\varepsilon(t,x)), & t\in\RR
\\[5pt]
X_\varepsilon(0,x)=x\in\RR^d,
\end{cases}
\end{equation}
and let $X$ be the limit flow associated with the limit vector field $a$.

Then, from any sequence of flows $X_\varepsilon$ we may derive a general sequence of vector fields $b_\varepsilon$ inducing the homogenization of the transport equation \eqref{eque}.
\begin{Pro}\label{pro.ae}
Let $u_\varepsilon^0$ be a bounded sequence in $L^p(\RR^N)$ with $p\in(1,\infty)$.
Assume that conditions \eqref{cae}, \eqref{aeA}, \eqref{divae} hold true. For a fixed $t>0$, define the vector field $W_\varepsilon:=X_\varepsilon(t,\cdot)$ from $\RR^N$ into $\RR^N$, and the vector field $b_\varepsilon$ by \eqref{siebe} with $\sigma_\varepsilon=1$. Then, the sequences $W_\varepsilon$ and $b_\varepsilon$ satisfy the assumptions of Theorem~\ref{thm.hom}.

Moreover, for any sequence $u_\varepsilon^0$ converging weakly to $u^0$ in $L^p(\RR^N)$, the solution $u_\varepsilon$ to equation~\eqref{eque} converges weakly in $L^\infty(0,T;L^p(\RR^N))$ to a solution $u$ to the equation \eqref{equ} where $\sigma_0=1$ and $\xi_0={\rm Cof}\,(D_x X(t,x))\,e_1$.
\end{Pro}
\begin{Rem}
There is a strong correspondence between the conditions \eqref{cae}-\eqref{aeA} and \eqref{divae} satisfied by the vector field $a_\varepsilon$, and respectively the conditions \eqref{cWe} and \eqref{cthe} satisfied by the vector fields $W_\varepsilon$ and $b_\varepsilon$.
\end{Rem}

\begin{proof}[Proof of Proposition~\ref{pro.ae}]
First of all, conditions \eqref{siebe} and \eqref{beB} are straightforward, since $\sigma_\varepsilon=1$ and $b_\varepsilon$ is divergence free.
Fix $T>0$. By \eqref{aeA} we have
\begin{equation}\label{eXe}
\forall\,t\in[0,T],\ \forall\,x\in\RR^N,\quad |X_\varepsilon(t,x)-x|\leq A\,T,
\end{equation}
so that the uniform property \eqref{upWe} is satisfied by $W_\varepsilon$.

Let $K$ be a compact set of $\RR^N$. Again by \eqref{eXe} there exists a compact subset $K'$ of~$\RR^N$ such that
\begin{equation}\label{K'}
\big\{X_\varepsilon(t,x),\,(t,x)\in[0,T]\times K\big\}\subset K'.
\end{equation}
Let $\de>0$. Since $a_\varepsilon$ converges uniformly to $a$ in $K'$ and $a\in C^1(\RR^N)$ is $k$-Lipschitz in~$K'$ for some $k>0$, we have for any small enough $\varepsilon>0$ and for any $t\in[0,T]$, for any $x,y\in K$,
\begin{align*}
\big|X_\varepsilon(t,x)-X_\varepsilon(t,y)\big| & \leq |x-y|+\int_0^t \big|a_\varepsilon(X_\varepsilon(s,x))-a_\varepsilon(X_\varepsilon(s,y))\big|\,ds
\\
& \leq \de+|x-y|+k\int_0^t\big|X_\varepsilon(s,x)-X_\varepsilon(s,y)\big|\,ds.
\end{align*}
Hence, by Gronwall's inequality (see, \eg \cite[\S 17.3]{HSD}) we get that for any small enough $\varepsilon>0$,
\[
\forall\,t\in[0,T],\ \forall\,x,y\in K,\quad |X_\varepsilon(t,x)-X_\varepsilon(t,y)\big|\leq(\de+|x-y|)\,e^{kt},
\]
which by \eqref{aeA} implies that for any small enough $\varepsilon>0$,
\[
\forall\,s,t\in[0,T],\ \forall\,x,y\in K,\quad |X_\varepsilon(s,x)-X_\varepsilon(t,y)\big|\leq A\,|s-t|+(\de+|x-y|)\,e^{kt},
\]
namely $X_\varepsilon$ is uniformly equicontinuous in the compact set $[0,T]\times K$. Therefore, by virtue of Ascoli's theorem this combined with \eqref{K'} and \eqref{cae} implies that up to a subsequence $X_\varepsilon$ converges uniformly in $[0,T]\times K$ to a solution $X$ to
\[
\forall\,t\in[0,T],\ \forall\,x\in K,\quad X(t,x)=x+\int_0^t a(X(s,x))\,ds,
\]
\ie $X$ is the flow associated with the vector field $a$. Since $a$ belongs to $C^1_b(\RR^N)$, the flow $X$ is uniquely determined (see, \eg \cite[\S 17.4]{HSD}).
Therefore, the whole sequence~$X_\varepsilon$ converges uniformly to $X$ in $[0,T]\times K$.
Moreover, by the differentiability of the flow (see, \eg \cite[\S 17.6]{HSD}) we have
\begin{equation}\label{DX}
\forall\,t\in[0,T],\ \forall\,x\in K,\quad D_xX_\varepsilon(t,x)=I_N+\int_0^t D_x X_\varepsilon(s,x)\,D_x a_\varepsilon(X_\varepsilon(s,x))\,ds,
\end{equation}
which using \eqref{cae}, \eqref{K'} and Gronwall's inequality implies that there exists a constant $c>0$ such that
\[
\forall\,t\in[0,T],\ \forall\,x\in K,\quad |D_x X_\varepsilon(t,x)|\leq |I_N|\,e^{ct}.
\]
Therefore, convergences \eqref{cWe} hold true.

On the other hand, by the Liouville formula associated with equation \eqref{DX} and estimate \eqref{aeA} we get that there exists a constant $c>1$ such that
\begin{multline}\label{edetXe}
\forall\,t\in[0,T],\ \forall\,x\in K,\\[-5pt]
c^{-1}\leq \det\left(D_xX_\varepsilon(t,x)\right)=\exp\left(\int_0^t (\div\,a_\varepsilon)(X_\varepsilon(s,x))\,ds\right)\leq c,
\end{multline}
which implies the existence of a constant $C>0$ such that for any $t\in[0,T]$ and $x\in K$,
\begin{multline*}
\big|\det\,(D_xX_\varepsilon(t,x))-\det\,(D_xX(t,x))\big|
\\
\leq C\int_0^T\!\!|\div\,a_\varepsilon-\div\,a|(X_\varepsilon(s,x))\,ds\\[-5pt]
+C\int_0^T\!\!\big|(\div\,a)(X_\varepsilon(s,x))-(\div\,a)(X(s,x))\big|\,ds.
\end{multline*}
Hence, using successively Jensen's inequality with respect to the integral in $s$, Fubini's theorem and the change of variables $y=X_\varepsilon(s,x)$ together with \eqref{K'} and \eqref{edetXe}, it follows that there exists a constant $C>0$ such that for any $t\in[0,T]$,
\begin{multline*}
\big\|\det\left(D_x X_\varepsilon(t,\cdot)\right)-\det\left(D_x X(t,\cdot)\right)\big\|_{L^q(K)}
\\
\leq C\,\|\div\,a_\varepsilon-\div\,a\|_{L^q(K')}+C\,\sup_{[0,T]\times K}\big|(\div\,a)(X_\varepsilon)-(\div\,a)(X)\big|.
\end{multline*}
This combined with convergence \eqref{divae} and the uniform convergence of $X_\varepsilon$ to $X$ in the compact set $[0,T]\times K$ implies the convergence \eqref{cthe} of $\theta_\varepsilon=\det (D_xX_\varepsilon(t,\cdot))$.

Finally, let $u_\varepsilon^0$ be a sequence in $L^p(\RR^N)$ converging weakly to $u^0$ in $L^p(\RR^N)$. By virtue of Theorem~\ref{thm.hom} combined with Remark~\ref{rem.vsiu} and recalling that $\sigma_\varepsilon=1$, the sequence $u_\varepsilon$ converges weakly in $L^p(\RR^N)$ to a solution $u$ to the equation \eqref{equ} where $\sigma_0=1$ and by \eqref{xi0cofDW}
\[
\xi_0={\rm Cof}\,(D_x X(t,\cdot))\,e_1\quad\mbox{in }\RR^N.
\]
Proposition~\ref{pro.ae} is thus proved.
\end{proof}

\backmatter
\bibliographystyle{jepplain+eid}
\bibliography{briane}
\end{document}
